\documentclass[12pt,a4paper]{amsart}

\usepackage[ngerman]{babel}
\usepackage{amsmath,amssymb,amsthm}
\usepackage{mathtools}
\usepackage[colorlinks=true,linkcolor=blue,citecolor=blue]{hyperref}
\usepackage{geometry}
\geometry{margin=2.5cm}

% -------------------------------------------------------
% Theorem-Umgebungen
% -------------------------------------------------------
\newtheorem{theorem}{Satz}[section]
\newtheorem{lemma}[theorem]{Lemma}
\newtheorem{korollar}[theorem]{Korollar}
\newtheorem{proposition}[theorem]{Proposition}
\theoremstyle{definition}
\newtheorem{definition}[theorem]{Definition}
\newtheorem{bemerkung}[theorem]{Bemerkung}

% -------------------------------------------------------
% Eigene Makros
% -------------------------------------------------------
\newcommand{\N}{\mathbb{N}}
\newcommand{\Z}{\mathbb{Z}}
\newcommand{\Q}{\mathbb{Q}}
\newcommand{\Fp}{\mathbb{F}_p}
\DeclareMathOperator{\ord}{ord}
\DeclareMathOperator{\ggT}{ggT}
\newcommand{\leg}[2]{\left(\frac{#1}{#2}\right)}

% -------------------------------------------------------
\begin{document}
% -------------------------------------------------------

\title{Anwendungen des Satzes von Wilson}

\author{Michael Fuhrmann}
\address{Specialist-Mathematics-Projekt, Build 66}
\email{specialist-maths@localhost}
\date{\today}

\begin{abstract}
Wir stellen vier wesentliche Anwendungen des Satzes von Wilson vor:
(i) Eine vollständige Primzahlcharakterisierung via $(n-1)! \equiv -1 \pmod{n}$
mit einer Analyse ihrer rechnerischen Unpraktikabilität;
(ii) einen Beweis, dass $-1$ genau dann ein quadratischer Rest modulo einer
ungeraden Primzahl $p$ ist, wenn $p \equiv 1 \pmod{4}$ gilt, direkt aus Wilson
via der Faktorisierung von $(p-1)!$ hergeleitet;
(iii) einen elementaren Beweis des harmonischen Lemmas von Wolstenholme:
$1 + 1/2 + \cdots + 1/(p-1) \equiv 0 \pmod{p}$ für alle Primzahlen $p \geq 3$;
und (iv) einen alternativen Beweis des kleinen Satzes von Fermat mit Hilfe
des Satzes von Wilson, mit einer Skizze, wie dieser Ansatz auf den Satz von
Euler erweitert werden kann.
Alle Beweise sind elementar und in sich abgeschlossen.
\end{abstract}

\maketitle
\tableofcontents

% ================================================================
\section{Einleitung}
% ================================================================

Der Satz von Wilson — $(p-1)! \equiv -1 \pmod{p}$ genau dann, wenn $p$
eine Primzahl ist — ist mehr als eine isolierte Kuriosität. Er ist ein
Werkzeug zur Herleitung tieferer Erkenntnisse über Primzahlen, quadratische
Reste und multiplikative Arithmetik.

In diesem Paper sammeln wir vier Anwendungen. Die erste ist der Satz selbst,
verwendet als Primalitätstest, mit einer ehrlichen Einschätzung seiner
Grenzen. Die zweite leitet ein vollständiges Kriterium dafür ab, dass $-1$
ein quadratischer Rest modulo $p$ ist — ein Ergebnis, das normalerweise via
Charaktertheorie oder Gauß-Summen bewiesen wird — hier mit nichts mehr als
Schulalgebra erhalten. Die dritte beweist eine harmonische Kongruenz von
Wolstenholme. Die vierte gibt einen alternativen Beweis des kleinen Satzes
von Fermat via Wilson und zeigt die Verflechtung dieser klassischen Ergebnisse.

% ================================================================
\section{Wilson als Primalitätstest}
% ================================================================

\begin{theorem}[Primzahlcharakterisierung]\label{satz:primalitaet}
  Eine ganze Zahl $n \geq 2$ ist genau dann eine Primzahl, wenn
  \[
    (n-1)! \;\equiv\; -1 \pmod{n}.
  \]
\end{theorem}

\begin{proof}
  \textbf{($\Rightarrow$)} Sei $n = p$ eine Primzahl. Die multiplikative
  Gruppe $(\Z/p\Z)^*$ hat Ordnung $p-1$ und ihre Elemente sind $1, 2, \ldots, p-1$.
  Jedes Element $a$ hat ein eindeutiges Inverses $a^{-1}$ in $\{1, \ldots, p-1\}$.
  Die einzigen Elemente, die gleich ihrem eigenen Inversen sind, erfüllen
  $a^2 \equiv 1 \pmod{p}$, d.\,h.\ $p \mid (a-1)(a+1)$; da $p$ prim ist,
  gilt $a \equiv \pm 1 \pmod{p}$, also $a \in \{1, p-1\}$.

  Alle anderen Elemente $2, \ldots, p-2$ bilden Paare $a \leftrightarrow a^{-1}$
  mit $a \cdot a^{-1} \equiv 1 \pmod{p}$. Daher:
  \[
    (p-1)! = 1 \cdot (p-1) \cdot \prod_{\substack{a=2\\a \neq a^{-1}}}^{p-2} a
    \equiv 1 \cdot (p-1) \cdot 1 \equiv p-1 \equiv -1 \pmod{p}.
  \]

  \textbf{($\Leftarrow$)} Sei $n$ zusammengesetzt, $n = ab$ mit $1 < a \leq b < n$.

  \emph{Fall 1: $a < b$.} Dann sind $a$ und $b$ verschiedene ganze Zahlen in
  $\{1, \ldots, n-1\}$, also gilt $ab = n \mid (n-1)!$ und
  $(n-1)! \equiv 0 \not\equiv -1$.

  \emph{Fall 2: $a = b$, also $n = a^2$.}
  Falls $a \geq 3$: Sowohl $a$ als auch $2a$ liegen in $\{1, \ldots, n-1\}$
  (da $2a < a^2$ für $a \geq 3$), also gilt $2a^2 = 2n \mid (n-1)!$ und
  $(n-1)! \equiv 0$.
  Falls $a = 2$: $n = 4$ und $(4-1)! = 6 \equiv 2 \pmod{4} \neq -1$.
\end{proof}

\begin{bemerkung}[Rechnerische Unpraktikabilität]
  Obwohl Satz~\ref{satz:primalitaet} theoretisch einen perfekten Primalitätstest
  liefert, ist er für große $n$ rechnerisch nutzlos. Die naïve Berechnung von
  $(n-1)! \bmod n$ erfordert $n-2$ Multiplikationen modulo $n$.
  Für $n \approx 2^{1000}$ (eine 300-stellige Primzahl) sind das $\sim 2^{1000}$
  Schritte — weit jenseits jeder praktikablen Berechnung.

  Im Gegensatz dazu arbeitet der Miller--Rabin-probabilistische Primalitätstest
  in $O(k \log^2 n)$ Zeit mit $k$ Runden, und der AKS-Algorithmus
  \cite{Agrawal2004} entscheidet die Primalität deterministisch in
  Polynomialzeit $O(\log^{6} n)$. Wilsons Test hat theoretische Bedeutung,
  aber keine praktische Anwendung für Primalitätstests.
\end{bemerkung}

% ================================================================
\section{Quadratische Reste: Wann ist $-1$ ein Quadrat modulo $p$?}
% ================================================================

\begin{definition}[Quadratischer Rest]\label{def:quadratischer_rest}
  Sei $p$ eine ungerade Primzahl. Eine ganze Zahl $a$ mit $\ggT(a,p) = 1$
  ist ein \emph{quadratischer Rest} modulo $p$, wenn $a \equiv x^2 \pmod{p}$
  für eine ganze Zahl $x$ gilt. Das Legendre-Symbol ist definiert durch
  \[
    \leg{a}{p} =
    \begin{cases}
      1  & \text{falls } a \text{ ein quadratischer Rest mod } p \text{ ist}, \\
      -1 & \text{falls } a \text{ ein Nicht-Rest mod } p \text{ ist.}
    \end{cases}
  \]
\end{definition}

\begin{theorem}[Erste Ergänzung zum quadratischen Reziprozitätsgesetz]%
  \label{satz:minus1_qr}
  Sei $p$ eine ungerade Primzahl. Dann ist $-1$ genau dann ein quadratischer
  Rest modulo $p$, wenn $p \equiv 1 \pmod{4}$ gilt. Äquivalent:
  \[
    \leg{-1}{p} \;=\; (-1)^{(p-1)/2}.
  \]
\end{theorem}

\begin{proof}
  Wir berechnen $(p-1)!$ auf zwei verschiedene Weisen.

  \textbf{Weise 1 (Wilson):}
  Nach dem Satz von Wilson gilt $(p-1)! \equiv -1 \pmod{p}$.

  \textbf{Weise 2 (Umgruppierung):}
  Schreibe das Produkt $(p-1)! = 1 \cdot 2 \cdots \frac{p-1}{2} \cdot
  \frac{p+1}{2} \cdots (p-1)$. In der zweiten Hälfte ersetze $j$ durch
  $p - j$ für $j \in \{(p+1)/2, \ldots, p-1\}$, was bijektiv auf
  $\{1, \ldots, (p-1)/2\}$ abbildet:
  Für $k = 1, 2, \ldots, (p-1)/2$ gilt $p - k \equiv -k \pmod{p}$.
  Also:
  \[
    (p-1)! = \prod_{k=1}^{(p-1)/2} k \cdot \prod_{k=1}^{(p-1)/2}(p-k)
    \equiv \prod_{k=1}^{(p-1)/2} k \cdot \prod_{k=1}^{(p-1)/2}(-k)
    = \Bigl(\prod_{k=1}^{(p-1)/2} k\Bigr)^2 \cdot (-1)^{(p-1)/2} \pmod{p}.
  \]
  Sei $M = \bigl(\tfrac{p-1}{2}\bigr)!$. Dann gilt:
  \[
    (p-1)! \equiv M^2 \cdot (-1)^{(p-1)/2} \pmod{p}.
  \]

  \textbf{Kombination:}
  Aus Weise 1 und Weise 2 folgt:
  \[
    -1 \equiv M^2 \cdot (-1)^{(p-1)/2} \pmod{p}.
  \]
  Daraus ergibt sich:
  \begin{itemize}
    \item Falls $(p-1)/2$ gerade ist (d.\,h.\ $p \equiv 1 \pmod{4}$):
      $(-1)^{(p-1)/2} = 1$, also $-1 \equiv M^2 \pmod{p}$.
      Damit ist $-1$ ein quadratischer Rest modulo $p$. \checkmark
    \item Falls $(p-1)/2$ ungerade ist (d.\,h.\ $p \equiv 3 \pmod{4}$):
      $(-1)^{(p-1)/2} = -1$, also $-1 \equiv -M^2 \pmod{p}$,
      d.\,h.\ $M^2 \equiv 1 \pmod{p}$.
      Wäre $-1$ ein quadratischer Rest, so gäbe es $x$ mit $-1 \equiv x^2 \pmod{p}$.
      Multiplikation liefert: $(-1)^{(p-1)/2} \equiv x^{p-1} \equiv 1 \pmod{p}$
      nach dem kleinen Satz von Fermat. Aber $(-1)^{(p-1)/2} = -1 \neq 1$.
      Widerspruch. Also ist $-1$ kein Rest. \checkmark
  \end{itemize}
  In beiden Fällen gilt $\leg{-1}{p} = (-1)^{(p-1)/2}$. \qedhere
\end{proof}

\begin{korollar}[Voraussetzung für Fermats Zwei-Quadrate-Satz]%
  \label{kor:zwei_quadrate}
  Eine Primzahl $p$ lässt sich als $p = x^2 + y^2$ mit ganzen Zahlen $x, y$
  schreiben, nur wenn $p = 2$ oder $p \equiv 1 \pmod{4}$.
\end{korollar}

\begin{proof}
  Falls $p = x^2 + y^2$, gilt $x^2 \equiv -y^2 \pmod{p}$. Falls $p \nmid y$,
  dann ist $(xy^{-1})^2 \equiv -1 \pmod{p}$, also ist $-1$ ein quadratischer
  Rest mod $p$, was $p \equiv 1 \pmod{4}$ nach Satz~\ref{satz:minus1_qr}
  erfordert. Der Fall $p = 2 = 1^2 + 1^2$ ist direkt. \qedhere
\end{proof}

% ================================================================
\section{Das harmonische Lemma von Wolstenholme}
% ================================================================

\begin{theorem}[Harmonische Kongruenz von Wolstenholme]\label{satz:harmonisch}
  Für jede Primzahl $p \geq 3$ gilt:
  \[
    1 + \frac{1}{2} + \frac{1}{3} + \cdots + \frac{1}{p-1}
    \;\equiv\; 0 \pmod{p},
  \]
  wobei jeder Bruch $1/j$ das multiplikative Inverse von $j$ modulo $p$ bezeichnet.
\end{theorem}

\begin{proof}
  Bezeichne die harmonische Summe mit $H = \sum_{j=1}^{p-1} j^{-1} \pmod{p}$.

  \textbf{Schritt 1 (Symmetrie).}
  Für jedes $j \in \{1, \ldots, p-1\}$ liegt das Inverse $j^{-1}$ ebenfalls
  in $\{1, \ldots, p-1\}$ (da $(\Z/p\Z)^*$ eine Gruppe ist), und die Abbildung
  $j \mapsto j^{-1}$ ist eine Bijektion auf $\{1, \ldots, p-1\}$.
  Daher ist die Summe über alle $j$ gleich der Summe über alle $j^{-1}$:
  \[
    H = \sum_{j=1}^{p-1} j^{-1} = \sum_{j=1}^{p-1} j \pmod{p}.
  \]

  \textbf{Schritt 2 (Auswertung von $\sum_{j=1}^{p-1} j$).}
  \[
    \sum_{j=1}^{p-1} j = \frac{(p-1)p}{2}.
  \]
  Modulo $p$: Da $p \equiv 0 \pmod{p}$, gilt $(p-1)p/2 \equiv 0 \pmod{p}$.

  \textbf{Schlussfolgerung.}
  Daher gilt $H \equiv \sum_{j=1}^{p-1} j \equiv 0 \pmod{p}$. \qedhere
\end{proof}

\begin{bemerkung}
  Wolstenholmes stärkerer Satz (1862) besagt, dass für $p \geq 5$ gilt:
  \[
    1 + \frac{1}{2} + \cdots + \frac{1}{p-1} \equiv 0 \pmod{p^2}.
  \]
  Dies ist ein tieferes Ergebnis und erfordert eine Analyse modulo $p^2$,
  die es mit dem Wilson-Quotienten verbindet (vgl.\ Paper~10 dieser Serie).
  Eine Primzahl $p$ ist genau dann eine Wilson-Primzahl, wenn die noch
  stärkere Kongruenz des Wilson-Quotienten modulo $p$ gilt.
\end{bemerkung}

% ================================================================
\section{Der kleine Satz von Fermat via Wilson}
% ================================================================

\begin{theorem}[Kleiner Satz von Fermat]\label{satz:fermat}
  Sei $p$ eine Primzahl und $a$ eine ganze Zahl mit $\ggT(a,p) = 1$.
  Dann gilt:
  \[
    a^{p-1} \;\equiv\; 1 \pmod{p}.
  \]
\end{theorem}

\begin{proof}[Beweis via Wilson]
  Betrachte die $p-1$ Vielfachen $a, 2a, 3a, \ldots, (p-1)a$.
  Wir behaupten, dass diese alle verschieden modulo $p$ sind und eine
  Permutation von $\{1, 2, \ldots, p-1\}$ bilden.

  \textbf{Schritt 1 (Verschiedenheit).}
  Falls $ja \equiv ka \pmod{p}$ mit $1 \leq j, k \leq p-1$, dann gilt
  $p \mid (j-k)a$. Da $\ggT(a,p)=1$, folgt $p \mid j-k$.
  Aber $|j-k| < p$, also $j = k$.

  \textbf{Schritt 2 (Sie liegen in $\{1,\ldots,p-1\}$).}
  Jedes $ja$ erfüllt $\ggT(ja, p) = \ggT(j,p)\ggT(a,p) = 1$
  (da $1 \leq j \leq p-1$ und $\ggT(a,p) = 1$), also gilt $ja \not\equiv 0 \pmod{p}$.

  \textbf{Schritt 3 (Produktvergleich).}
  Wir multiplizieren alle $p-1$ Kongruenzen:
  \[
    a \cdot 2a \cdot 3a \cdots (p-1)a \;\equiv\; 1 \cdot 2 \cdot 3 \cdots (p-1) \pmod{p}.
  \]
  Die linke Seite ist $a^{p-1} \cdot (p-1)!$ und die rechte Seite ist $(p-1)!$.
  Daher gilt:
  \[
    a^{p-1} \cdot (p-1)! \equiv (p-1)! \pmod{p}.
  \]

  \textbf{Schritt 4 (Kürzung via Wilson).}
  Nach dem Satz von Wilson gilt $(p-1)! \equiv -1 \pmod{p}$, also ist
  $\ggT((p-1)!, p) = 1$ (da $p \nmid (p-1)!$ klar ist, weil $p$ prim und
  keines der Elemente $1, \ldots, p-1$ durch $p$ teilbar ist).
  Daher ist $(p-1)!$ invertierbar modulo $p$, und wir können kürzen:
  \[
    a^{p-1} \equiv 1 \pmod{p}. \qedhere
  \]
\end{proof}

\begin{bemerkung}
  Der Standardbeweis des kleinen Satzes von Fermat (der das Permutationsargument
  aus Schritten 1--3 verwendet) benötigt Wilson eigentlich nicht — das Kürzen
  in Schritt 4 erfordert nur $(p-1)! \not\equiv 0 \pmod{p}$, was offensichtlich
  ist, da $p$ prim und keines der Elemente $1, \ldots, p-1$ durch $p$ teilbar ist.
  Der Satz von Wilson wird hier verwendet, um den \emph{Wert} von $(p-1)!$ zu
  liefern und den Kürzungsschritt expliziter darzustellen.
  Der obige Beweis zeigt, wie Wilson und Fermat Facetten derselben
  gruppentheoretischen Tatsache sind: $(\Z/p\Z)^*$ ist eine Gruppe der Ordnung $p-1$.
\end{bemerkung}

\begin{korollar}[Skizze: Eulerscher Satz via Wilson + Struktur]\label{kor:euler}
  Für $\ggT(a, n) = 1$ gilt: $a^{\varphi(n)} \equiv 1 \pmod{n}$.
\end{korollar}

\begin{proof}[Beweisskizze]
  Nach dem Chinesischen Restsatz gilt $(\Z/n\Z)^* \cong \prod_i (\Z/p_i^{k_i}\Z)^*$
  für die Primfaktorzerlegung $n = \prod p_i^{k_i}$.
  Jeder Faktor $(\Z/p_i^{k_i}\Z)^*$ hat Ordnung $\varphi(p_i^{k_i})$, also gilt
  $a^{\varphi(p_i^{k_i})} \equiv 1 \pmod{p_i^{k_i}}$.
  Da $\varphi(n) = \prod \varphi(p_i^{k_i})$ und $\varphi(p_i^{k_i}) \mid \varphi(n)$,
  erhalten wir $a^{\varphi(n)} \equiv 1$ modulo jedem $p_i^{k_i}$, also
  $a^{\varphi(n)} \equiv 1 \pmod{n}$ nach CRT.
  Die Rolle Wilson-artiger Sätze ist die Sicherstellung, dass die Einheitengruppen
  $(\Z/p^k\Z)^*$ die Ordnung genau $\varphi(p^k) = p^{k-1}(p-1)$ haben, was die
  Zyklizität dieser Gruppen erfordert (bewiesen mit der Existenz von Primitivwurzeln,
  die wiederum Wilson für $p$ verwendet).
\end{proof}

% ================================================================
\section{Schlussfolgerung}
% ================================================================

Der Satz von Wilson ist zwar als Primalitätstest rechnerisch unpraktisch,
beleuchtet aber tiefe strukturelle Eigenschaften der ganzen Zahlen modulo $p$.

Die Herleitung der ersten Ergänzung zum quadratischen Reziprozitätsgesetz
($\leg{-1}{p} = (-1)^{(p-1)/2}$) aus Wilson durch eine doppelte Zählung
von $(p-1)!$ ist eines der elegantesten Argumente in der elementaren
Zahlentheorie. Das harmonische Lemma von Wolstenholme folgt aus einem einfachen
Symmetrieargument über die Gruppe $(\Z/p\Z)^*$. Und der Beweis des kleinen
Satzes von Fermat via Wilson zeigt die gruppentheoretische Einheit beider
Ergebnisse.

Diese vier Anwendungen zeigen, dass der Satz von Wilson keine isolierte
Kuriosität ist, sondern ein grundlegendes Werkzeug, das Primalität,
quadratische Reziprozität, harmonische Summen und die multiplikative Struktur
von $\Z/p\Z$ verbindet.

\begin{thebibliography}{9}

\bibitem{Hardy1979}
G.~H.~Hardy und E.~M.~Wright.
\newblock \emph{Einführung in die Zahlentheorie} (An Introduction to the Theory of Numbers), 5.~Aufl.
\newblock Oxford University Press, 1979.

\bibitem{Ireland1990}
K.~Ireland und M.~Rosen.
\newblock \emph{A Classical Introduction to Modern Number Theory}, 2.~Aufl.
\newblock Springer, New York, 1990.

\bibitem{Agrawal2004}
M.~Agrawal, N.~Kayal und N.~Saxena.
\newblock PRIMES is in P.
\newblock \emph{Ann.\ of Math.}, 160(2):781--793, 2004.

\bibitem{Niven1991}
I.~Niven, H.~S.~Zuckerman und H.~L.~Montgomery.
\newblock \emph{An Introduction to the Theory of Numbers}, 5.~Aufl.
\newblock Wiley, New York, 1991.

\bibitem{Crandall1997}
R.~Crandall, K.~Dilcher und C.~Pomerance.
\newblock A search for Wieferich and Wilson primes.
\newblock \emph{Math.\ Comp.}, 66(217):433--449, 1997.

\end{thebibliography}

\end{document}
